%%=============================================================================
%% 第 5 章 总结与展望
%%=============================================================================

\chapter{总结与展望}
\label{chap:conclusion}

\section{研究总结}
\label{sec:summary}

运动想象脑机接口（Motor Imagery Brain-Computer Interface, MI-BCI）作为非侵入式脑机接口的重要分支，在康复医疗和人机交互等领域展现出广阔的应用前景。然而，时间窗选择作为影响特征提取质量和分类性能的关键因素，长期以来制约着 MI-BCI 系统的性能提升。现有方法（如固定时间窗、贝叶斯优化、深度强化学习）在面对数据稀缺性和设备资源约束时存在明显局限：固定时间窗无法适应个体间差异，贝叶斯优化计算复杂度高，深度强化学习方法模型复杂、训练不稳定且可解释性差。

针对上述问题，本研究提出一种基于表格型 Q-Learning 的轻量化时间窗优化方法。该方法将时间窗优化问题形式化为一个包含 136 个离散状态的马尔可夫决策过程（MDP），通过经典的时序差分学习迭代优化时间窗选择策略。在 BCI Competition IV 2a 数据集上的系统实验表明，本研究方法在分类性能上与深度 Q 网络（DQN）无显著差异（62.68\% vs 62.70\%, $p = 0.981$），同时显著优于固定时间窗方法（约 45\%）。更重要的是，本研究方法实现了 3 倍训练加速、1000 倍推理延迟降低和 125 倍内存占用减少，完美契合便携式 BCI 设备的资源约束需求。

本研究的主要贡献可概括为以下三个方面：

（1）\textbf{轻量化 BCI 设计新范式}。针对便携式 BCI 设备的资源约束，本研究提出将深度强化学习简化为表格型强化学习的设计思路，成功实现了轻量化与高性能的统一。这一设计范式为资源受限场景下的 BCI 系统优化提供了新思路。

（2）\textbf{算法选择准则的理论贡献}。本研究通过理论分析和实验验证，明确了表格型 Q-Learning 与 DQN 的适用边界：当状态空间规模小于$10^4$且状态 - 动作对可充分访问时，表格型方法在收敛性保证、样本效率和可解释性方面均优于 DQN。这一准则可推广至其他状态空间有限的 BCI 优化任务。

（3）\textbf{可解释性 BCI 系统的实践探索}。通过提供可直接可视化的 Q 值表，本研究增强了 BCI 系统在医疗康复领域应用的可解释性。在安全敏感应用中，可解释性是系统部署的关键前提，本研究为此提供了实践参考。

\section{研究局限性与不足}
\label{sec:limitations}

尽管本研究取得了预期的成果，但仍存在以下局限性：

（1）\textbf{验证范围的局限}。所有实验均基于 BCI Competition IV 2a 这一标准数据集上的健康被试数据，未能在真实患者群体（如脑卒中康复患者）或在线闭环 BCI 系统中进行验证。患者脑电信号的非平稳性更强、信噪比更低，可能对本方法的鲁棒性提出更大挑战。此外，数据集仅包含 9 名被试，样本规模有限，统计检验的效力可能不足。

（2）\textbf{方法设计的局限}。本研究对时间窗参数的离散化步长为 0.2 秒，这种近似处理可能无法定位到全局最优的连续时间点。尽管 0.2 秒步长在实验中表现良好，但更精细的离散化（如 0.1 秒）或连续优化方法可能带来性能提升。此外，本研究的奖励设计存在一定的局限性：时间窗优化的奖励信号只能通过间接辅助任务（即 CSP-SVM 分类准确率）来构建，而非直接衡量时间窗本身的质量。这种间接奖励设计可能导致优化目标偏差——分类准确率不仅受时间窗影响，还受 CSP 特征提取质量、SVM 参数设置等因素干扰。理想的奖励函数应能直接反映时间窗与 ERD 信号质量的关联\cite{pfurtscheller1999event}，但如何设计这样的奖励函数仍是一个开放性难题。未来可探索多源奖励信号，如结合 ERD 强度、信噪比估计、或基于信息论的度量（如互信息），以提供更丰富、更直接的优化指导。

（3）\textbf{Q-Learning 变体研究的不足}。本研究虽在消融实验中对比了 Standard Q-Learning、Double Q-Learning、Dueling Q-Learning 和 Dueling Double Q-Learning 四种变体，但研究深度仍有不足。首先，实验仅观察到 Double Q-Learning 在本任务中性能略低于 Standard Q-Learning 的现象，但未深入分析其内在机制——可能是由于双 Q 网络的过估计修正策略在状态空间有限、奖励确定性高的任务中过于保守，导致探索效率下降。其次，本研究未探索其他 Q-Learning 改进变体，如 Prioritized Experience Replay\cite{schaul2015prioritized}、Noisy Nets\cite{fortunato2018noisy}、Distributional RL\cite{bellemare2017distributional} 等，这些方法在提升样本效率或稳定性方面可能有潜在价值。最后，本研究未对 Q 值表的演化过程进行深入分析，如不同状态 - 动作对的收敛速度、Q 值分布的时空模式等，这些分析可能为理解表格型方法在 BCI 优化任务中的学习动态提供新的洞察。

（4）\textbf{通用性的局限}。本研究方法目前针对特定数据集（BCI Competition IV 2a）和 CSP-SVM 分类器设计，其在不同分类模型（如 EEGNet\cite{lawhern2018eegnet}、Shallow FBCSP\cite{ang2008filter} 等深度学习模型）或其他 BCI 范式（如 SSVEP、P300）上的有效性尚待验证。方法的泛化能力需要在更广泛的实验设置中进行评估。

（5）\textbf{分析维度的局限}。本研究对计算效率的分析主要停留在参数计数和理论复杂度层面，缺乏在 ARM 架构单片机等典型嵌入式 BCI 平台上的实际运行时间和功耗的基准测试。理论分析虽能说明趋势，但实际部署性能仍需实证验证。

\section{未来工作展望}
\label{sec:future_work}

基于本研究的成果与局限，未来工作将从以下几个方向展开：

（1）\textbf{开展在线和临床验证}。我们将把训练好的模型移植到基于嵌入式处理器的便携式 BCI 系统中，开展在线实验，评估其在实时反馈场景下的响应速度和用户体验。同时，积极寻求与康复机构的合作，在运动功能障碍患者群体中检验本方法的临床有效性和可用性。真实患者数据的验证将为本方法的实际应用价值提供更有力的证据。

（2）\textbf{进一步优化算法效率}。针对离线训练成本高的问题，我们将研究基于元学习的快速奖励估计方法，或采用缓存机制复用交叉验证结果，以期将离线训练时间从分钟级缩短到秒级，真正实现\"即插即用\"的用户校准。此外，探索更精细的离散化步长（如 0.1 秒）或结合连续优化方法（如贝叶斯优化）的可能性，以进一步提升时间窗选择精度。

（3）\textbf{探索算法的泛化能力}。我们将尝试将本方法与其他特征提取器（如 EEGNet\cite{lawhern2018eegnet}）和分类器（如深度学习模型）结合，构建模块化的自适应 BCI 系统，验证本方法在不同技术栈下的有效性。更重要的是，我们将把该表格型强化学习框架推广至运动想象的\textbf{空间层面优化}任务中，包括：（a）\textbf{通道选择}：基于强化学习选择最优 EEG 通道子集，减少冗余通道、降低计算负担\cite{mcfarland1997spatial}；（b）\textbf{CSP 空间滤波器优化}：联合优化 CSP 分量数和空间滤波器参数，提升特征判别力\cite{lotte2011regularizing}；（c）\textbf{时空联合优化}：参考 Zhang 等人（2019）的时间约束稀疏组空间模式（TSGSP）思路\cite{zhang2019temporally}，将时间窗选择与空间滤波优化统一建模为马尔可夫决策过程，实现时空特征的协同优化。通过上述扩展，验证表格型强化学习作为通用轻量化自适应工具的潜力。

（4）\textbf{研究迁移学习策略}。计划对不同被试的最优 Q 表进行分析，尝试从中提取通用的时间窗选择先验知识，从而在新用户仅提供极少量试次（如 10--20 个试次）的情况下，快速迁移并微调出一个个性化策略，大幅缩短系统校准时间。迁移学习策略的实现将使本方法更接近实际应用场景，降低新用户的使用门槛。

